\chapter{Background and related work}
\label{ch:background_related}

\section{Brown bear ecology and demographic context}

Brown bear ecology is a very good case for simulation, because the population trends we read about in reports are never the product of a single cause. They emerge from many local processes acting at the same time: reproduction, mortality at different ages, access to food, the way individuals use space, and the constant friction of living next to humans. In practice none of these processes acts alone. They overlap in time and they interact, and their combined effect is what pushes a population toward growth, decline, or a local equilibrium. This is precisely why a verbal, one-cause-at-a-time explanation tends to fall short, and why a model that lets the causes act together is useful.

From the demographic literature we know that brown bear trajectories are shaped above all by age and sex structure, by the long intervals between successive litters, by cub survival, and by the mortality pressure that human-dominated landscapes add on top of natural mortality \cite{Schwartz2003,Saether1998,Swenson1997}. These are not abstract details for a modeller; they are the realistic bounds that keep the model honest. When the simulation in this thesis uses a minimum reproductive age of four years, a gestation of about a year, a multi-year cooldown between litters, a litter mean close to two cubs, and elevated mortality for cubs, none of those numbers is invented. They come from this body of work, and the role of the literature is exactly to provide defensible ranges instead of arbitrary values \cite{Swenson2000}.

The Romanian context makes the problem more relevant still. Romania hosts by a wide margin the largest brown bear population in the European Union, and recent work has documented both the movement ecology of these bears and the patterns of human--bear interaction in the Carpathians \cite{Pop2018,Popescu2017}. At the same time, institutional reports and policy documents point to substantial changes both in the estimated size of the population and in the management strategy applied to it over the last two decades \cite{INCDS2023,MMAP2024}. The increase in conflict reports and the shifting estimates of population size are not just background colour; they are the very signal this thesis tries to explain. So this is not only a theoretical modelling exercise. It is a problem with direct public-policy weight, where the question of which mechanism actually drove the rise has practical consequences for how the population should be managed \cite{Chapron2014}.

\section{Simulation and the agent-based paradigm}

Agent-based modelling is the natural tool to reach for when individual heterogeneity matters and when interactions are local in space and time, and both of those are true for bears. Instead of forcing the whole population into one aggregated equation, an agent-based model lets each individual carry its own internal state and update its behaviour according to its own local context. The population-level behaviour is then not assumed in advance; it appears as an emergent consequence of all those individual updates. For a question like ours, where we want to test whether a given assumption is \emph{sufficient} to produce an observed trend, this emergence is the whole value of the method.

The methodological foundations of individual-based ecological models are by now well established \cite{GrimmRailsback2005,RailsbackGrimm2019,DeAngelis2005}. Just as importantly, the field has developed shared standards for describing such models, most notably the ODD protocol and its update, which lay out a practical structure for stating a model's entities, processes, scheduling and reproducibility decisions \cite{Grimm2006ODD,Grimm2020ODD}. This thesis follows the same spirit, even where it does not follow the protocol to the letter: Chapter~3 explicitly defines the agent state, the scheduling of the tick loop, the environment representation and the experiment protocol, precisely so that the model can be understood and, in principle, reproduced.

The reason this paradigm fits counterfactual analysis so well is that interventions can be expressed as explicit rules and then switched on and off. If we want to test the effect of changing hunting pressure or of adding a food subsidy, we do not have to re-derive an equation. We change a rule or a parameter at a given year and re-run, and because the model is stochastic, we re-run it several times and look at the distribution of outcomes rather than a single line. This is the same logic used across the wider agent-based modelling community, from general treatments of the method \cite{Macal2016} to teaching-oriented platforms such as NetLogo \cite{WilenskyRand2015}.

\section{Romanian case study framing}

In this thesis the Romanian situation is treated as a concrete testbed, not as a generic backdrop. More precisely, the empirical framing is the brown bear population across Romania's Carpathian forests as a whole, with 2005--2021 as the main analysis window. Choosing a real framing rather than a purely abstract one brings two clear advantages.

First, it provides real-world constraints. Demographic assumptions, habitat structure and the policy timeline can all be tied to published or institutional sources, which sharply reduces the number of arbitrary modelling choices and makes the scenarios defensible. The landscape transitions of the Romanian Carpathians over recent decades have themselves been studied \cite{Munteanu2014,Griffiths2013}, which gives some grounding to the habitat-loss scenario we test later. Second, a real framing allows the results to be interpreted in a management-oriented direction. Even though the model stays simplified, its scenario outputs can still be read as statements like ``this combination of pressures reproduces the observed trend more plausibly than that one'', which is the kind of statement a manager can actually use.

At the same time, this framing brings a responsibility in how the results are read. Institutional population estimates evolve as census methods improve, and policy periods do not always begin and end on clean calendar boundaries. For that reason the simulation is presented as a structured comparative tool under explicit assumptions, and not as an exact forecasting device. I would rather make a careful comparative claim that holds than a precise numerical claim that the data cannot support.

\section{Demography and the problem of parameter realism}

There is a recurring weakness in student simulation projects that I was determined to avoid, and it is worth naming directly: parameter arbitrariness. In ecological agent-based modelling it is dangerously easy to produce plots that look plausible but that mean very little, because the parameters behind them were tuned by hand until the picture looked nice. A model like that can be made to show almost anything, which means it shows nothing.

The defence against this is to constrain parameter ranges with the demographic literature wherever possible, and only then to calibrate within those ranges. The references used here support realistic intervals for the quantities that matter most: age-related mortality, the vulnerability and survival bottleneck of cubs, the intervals between litters, and the additional mortality that human pressure imposes \cite{Schwartz2003,Saether1998,McLellan1989,Bischof2009}. The sex- and age-specific mortality multipliers in the model, the elevated cub danger, and the long reproductive cooldown all sit inside ranges drawn from this work. This does not make the uncertainty disappear, but it does give a defensible starting point for the calibration and sensitivity experiments, which is a very different situation from picking numbers until the curve fits.

\section{Methodological references for evaluation}

Besides the ecological literature, the thesis also leans on methodological sources, because how you evaluate a stochastic simulation is itself a methodological question. Sensitivity analysis and parameter-estimation workflows for this kind of model are well covered in the standard references \cite{Saltelli2008,Thiele2014}, and they inform the way Chapter~4 reports its results. The decision that follows from them is simple but non-negotiable: if a claim depends on a stochastic simulation, a single run is not evidence. Every scenario in this thesis is therefore run as several seeded replicates and then aggregated, and the comparison against the reference trajectory is made on those aggregates rather than on one lucky or unlucky line.

\section{Related software approaches}

From the software side there is a mature body of prior work on multi-agent simulation environments, including MASON and Repast and their high-performance, distributed variants \cite{Luke2005,CollierNorth2013,North2013}, lighter frameworks such as Mesa \cite{Masad2015}, and dedicated spatially explicit ecological platforms \cite{Bocedi2014}. There is also a long-standing theoretical foundation for parallel and distributed simulation in general \cite{Fujimoto2000}. I chose to build the model directly in Java rather than on top of one of these frameworks, partly for full control over the execution and reproducibility machinery, and partly because the computational question I wanted to ask is about that very machinery.

What this literature makes clear, and what is central to my argument, is that the benefit of parallel or distributed execution depends entirely on the shape of the workload: on how heavy the per-agent computation is, on the communication pattern, and on how often the model has to synchronise. In models where each agent does heavy local computation between synchronisations, distribution scales well. In models where the per-agent work is light and synchronisation is frequent, communication overhead can erode the speedup and even reverse it. This distinction is exactly the one that applies to the bear model. Because it is tick-based and globally synchronised, and because each agent action is cheap, distributing a single run can be slower than running it multithreaded on one machine, while distributing many independent runs remains very effective. The wider ecological agent-based literature, including work on other large carnivores such as tigers \cite{Carter2015}, confirms more generally that complex wildlife dynamics can be represented faithfully at the individual level and produce meaningful scenario outputs, which is the methodological reassurance this thesis needed before committing to the approach.

\section{Gap addressed by this thesis}

The gap this thesis addresses has two sides, one ecological and one computational, and the contribution is to close both at once with a single piece of work.

On the ecological side, much of the public and even technical discussion about bear-population change in Romania stays descriptive, and it tends to emphasise one driver at a time, it was the hunting ban, or it was the feeding, or it was habitat change, without ever testing how these drivers behave together inside one consistent framework. This thesis addresses that by integrating several candidate drivers into one model and evaluating them, both in isolation and combined, under controlled scenarios with replicated runs.

On the computational side, reports and student projects alike often present multithreaded and distributed execution as if they were universally beneficial, almost as a synonym for ``more serious''. In practice, as the simulation-systems literature warns, that is not true for every workload. For this tick-synchronised model the thesis works out explicitly where parallelism actually helps, and, just as importantly, where it would not. The model is therefore executed with a multithreaded, single-machine design, and a distributed execution is deliberately \emph{not} built, because the analysis shows it would not pay off for this granularity. The reasoned conclusion is a simple two-level picture:

\begin{itemize}
    \item multithreading is used to accelerate a single simulation run on one machine, where the threads share memory and the per-tick synchronisation barrier is cheap;
    \item distribution would only make sense at the level of large batches of independent runs, which are embarrassingly parallel, but at the scale of these experiments a single multi-core machine ran the whole campaign comfortably, so it was not needed and not implemented.
\end{itemize}

Stating this judgement explicitly, rather than reflexively reaching for a distributed implementation, improves both the scientific reproducibility of the results and the engineering honesty of how they were produced, and it is one of the things that I believe distinguishes this work from a model that simply runs and prints a curve.

\section{Chapter summary}

This chapter set out the conceptual and bibliographic foundation of the thesis. The ecological context motivates the modelling problem and supplies realistic parameter ranges; the agent-based literature supports the choice of paradigm and the standards for describing and evaluating such a model; and the simulation-systems literature supplies the criteria for the computational part of the evaluation. With this foundation in place, the next chapter moves from background to implementation: how the world is represented, how agents decide and act, how their effects are committed each tick, and how the whole thing is wrapped in a reproducible experiment pipeline.
